% The official journal class is included in the source archive. The amsart
% fallback supports the standalone editor preview; release builds require baustms.
\newif\ifbulletinformat
\IfFileExists{baustms.cls}{\bulletinformattrue\documentclass[plain]{baustms}}{\bulletinformatfalse\documentclass[11pt]{amsart}}
\ifbulletinformat
  \usepackage{microtype,mathtools}
  \theoremstyle{cupthm}
  % Suppress the template's unassigned publication-DOI placeholder.
  \makeatletter
  \def\@mainheadline{\gdef\@mymainhead{}\setcounter{page}{1}}
  \def\ps@austmstitle{\let\@mkboth\@gobbletwo
    \def\@oddhead{}\def\@evenhead{}\def\@oddfoot{}\def\@evenfoot{}}
  \makeatother
\else
  \usepackage[T1]{fontenc}
  \usepackage{lmodern,microtype,mathtools,amssymb,hyperref}
\fi
\hypersetup{draft=false,hidelinks,pdftitle={A counterexample to Kourovka Notebook Problem 16.45},pdfauthor={Alec Kriebel}}
\newtheorem{theorem}{Theorem}
\newtheorem{proposition}[theorem]{Proposition}
\newtheorem{lemma}[theorem]{Lemma}
\ifbulletinformat\theoremstyle{cuprem}\else\theoremstyle{remark}\fi
\newtheorem{remark}[theorem]{Remark}
\newcommand{\F}{\mathbf F}
\newcommand{\mup}{\mu'}
\newcommand{\bfth}{b_{\mathrm f}}
\newcommand{\SL}{\operatorname{SL}}
\newcommand{\Sym}{\operatorname{Sym}}
\newcommand{\Core}{\operatorname{Core}}
\title{A counterexample to Kourovka Notebook Problem 16.45}
\ifbulletinformat
  \runningtitle{Kourovka Notebook Problem 16.45}
  \author[1]{Alec Kriebel}
  \address[1]{Independent researcher, San Francisco, USA\\
    ORCID: \href{https://orcid.org/0009-0001-9320-500X}{0009-0001-9320-500X}\email{me@aleckriebel.com}}
  \authorheadline{A. Kriebel}
  \classification[2020]{Primary 20B05; Secondary 20D30, 20D60}
\else
  \author{Alec Kriebel}
  \address{Independent researcher, San Francisco, USA}
  \email{me@aleckriebel.com}
  \urladdr{https://orcid.org/0009-0001-9320-500X}
  \subjclass[2020]{Primary 20B05; Secondary 20D30, 20D60}
\fi
\keywords{Permutation bases, independent sets, finite groups, subgroup lattices, Kourovka Notebook}
\date{}
\begin{document}
\begin{abstract}
Let $b(G)$ be the largest size of an inclusion-minimal base over all
permutation representations of a finite group $G$, with bases taken for the
permutation image, and let $\mup(G)$ be the largest size of an independent
subset of $G$. We construct an explicit affine group of order $100920$ for
which $b(G)=3<4=\mup(G)$, giving a negative answer to Kourovka Notebook
Problem 16.45. The maximum remains three when restricted to faithful
representations. The proof is structural: every subgroup missing the
translation subgroup has faithful base invariant at most two, and every
nontrivial normal subgroup contains the translations. The independent
four-set gives an irredundant intersection family whose intersection is
central in the complement but not normal in the affine group.
\end{abstract}
\maketitle

\section{Introduction}
A subset $S$ of a finite group $X$ is \emph{independent} if
$x\notin\langle S\setminus\{x\}\rangle$ for every $x\in S$.
Write $\mup(X)$ for the largest size of such a subset; it need not generate
$X$. For an action $\rho:X\to\Sym(\Omega)$, a \emph{base} is a set of
points with pointwise stabilizer $\ker\rho$. A base is \emph{minimal} if
it is inclusion-minimal. Let $b(X)$ be the largest size of a minimal base
over all permutation representations of $X$, and let $\bfth(X)$ denote
the same maximum restricted to faithful representations. In particular,
$b(1)=\bfth(1)=\mup(1)=0$. The definitions allow intransitive and
nonfaithful actions. Allowing infinite $\Omega$ makes no difference:
choosing a moved point for each nonidentity element of the finite image
gives a finite subbase of every base.

Cameron's Problem 16.45 in the \emph{Kourovka Notebook} asks whether
$b(X)=\mup(X)$ for every finite group $X$ \cite{notebook}.
The inequality $b(X)\le\mup(X)$ is known. Cameron uses $b_2$ for $b$,
and uses $\mu'$ in \cite{cameron2014} and $\mu^\uparrow$ in
\cite{cameron2024} for our independence invariant. In contrast,
Anagnostopoulou-Merkouri and Burness \cite[Remark 1(c)]{ab2024} restrict
their $b_2$ to faithful actions, so it corresponds to our $\bfth$.
The all-actions question is still posed in Cameron's 2024 account
\cite[Section 6, Proposition 2 and the following discussion]{cameron2024}
and is listed without a solution comment in the September 2026 notebook.
We prove the following.

\begin{theorem}\label{thm:main}
Over $\F_{29}$, put
\[
 A=\begin{pmatrix}0&28\\1&0\end{pmatrix},\qquad
 B=\begin{pmatrix}2&7\\12&28\end{pmatrix},\qquad
 H=\langle A,B\rangle\le\SL_2(29).
\]
For the natural affine semidirect product $G=\F_{29}^2\rtimes H$,
\[
 |G|=100920,\qquad \bfth(G)=b(G)=3<4=\mup(G).
\]
\end{theorem}

The proof separates two intersection conditions in the subgroup lattice.
Independence corresponds to irredundant intersection families with arbitrary
intersection, whereas permutation bases require that intersection to be
normal. In our example an independent four-set gives a family with
intersection of order two, central in $H$ but not normal in $G$.
The faithful upper bound follows by restricting an intersection family to
a subgroup that does not contain the translations. The normal-subgroup
structure then handles every nonfaithful action.
No enumeration of the subgroup lattice of $G$ is required, and no
minimum-order assertion is made.

\section{Subgroup formulations}
We recall Cameron's subgroup formulations
\cite[Propositions 3.1--3.2 and Corollary 3.3]{cameron2014};
see also \cite{cameron2024}. We include the short arguments to make the
normality condition explicit.

\begin{proposition}\label{prop:bases}
The invariant $b(X)$ is the largest $t$ for which subgroups
$L_1,\ldots,L_t$ of $X$ satisfy
\[
 N=\bigcap_{i=1}^t L_i\lhd X,\qquad
 \bigcap_{j\ne i}L_j>N\quad(1\le i\le t).
\]
The same description with $N=1$ gives $\bfth(X)$. Consequently,
\begin{equation}\label{eq:quotients}
 b(X)=\max_{N\lhd X}\bfth(X/N).
\end{equation}
\end{proposition}
\begin{proof}
Point stabilizers of a minimal base have the stated properties.
Conversely, let $X$ act on the disjoint union of the left coset spaces
$X/L_i$. The action kernel is
\[
 \bigcap_i\Core_X(L_i)
 =\Core_X\!\left(\bigcap_i L_i\right)=N.
\]
The identity cosets form a minimal base. Factoring through the action kernel
gives \eqref{eq:quotients}. Empty intersections are taken in the ambient
group, so the empty family also handles the trivial group and trivial actions.
\end{proof}

Call a subgroup family \emph{meet-irredundant} if every member is essential
to its total intersection, without a normality condition.

\begin{lemma}\label{lem:breadth}
The largest size of a meet-irredundant family in $X$ is $\mup(X)$.
In particular, $b(X)\le\mup(X)$.
\end{lemma}
\begin{proof}
For a meet-irredundant family $L_1,\ldots,L_t$, choose
$x_i\in(\bigcap_{j\ne i}L_j)\setminus L_i$.
Every $x_j$ with $j\ne i$ lies in $L_i$, whereas $x_i$ does not, so
the $x_i$ are distinct and independent. Conversely, for an independent
set $\{x_1,\ldots,x_t\}$, the subgroups
$L_i=\langle x_j:j\ne i\rangle$ are meet-irredundant, with witnesses $x_i$.
\end{proof}

\begin{lemma}[Restriction]\label{lem:restriction}
If $L_1,\ldots,L_t$ is meet-irredundant in $X$, then
$t-1\le\mup(L_i)$ for each $i$. If its total intersection is $1$,
then $t-1\le\bfth(L_i)$.
\end{lemma}
\begin{proof}
Inside $L_i$, the family $L_i\cap L_j$ for $j\ne i$ is meet-irredundant:
the witness for $L_j$ belongs to $L_i$. Its total intersection is unchanged.
Apply Lemma~\ref{lem:breadth} or Proposition~\ref{prop:bases}.
\end{proof}

These descriptions also show that $\mup$ and $\bfth$ are monotone under
subgroups. For completeness, a meet-irredundant family gives the
top-preserving Boolean meet-semilattice embedding
\[
 J\longmapsto\bigcap_{i\notin J}L_i
 \qquad(J\subseteq\{1,\ldots,t\}).
\]
The witnesses distinguish the images, and the bottom is $\bigcap_iL_i$.
Thus the normality distinction is exactly the one in the notebook's
semilattice formulation; no preservation of joins is assumed.

\section{The complement}
Let $H$ be as in Theorem~\ref{thm:main}, and set $Z=\{I,-I\}$.
Direct multiplication gives
\begin{equation}\label{eq:triangle}
 A^2=B^3=(AB)^5=-I.
\end{equation}
The spherical triangle presentation
$A_5=\langle a,b\mid a^2=b^3=(ab)^5=1\rangle$
and simplicity of $A_5$ show that $H/Z\cong A_5$: the quotient is
nontrivial because $A\notin Z$. Hence $|H|=120$.

We recall the elementary subgroup facts that will be used. Every proper
subgroup of $A_5$ lies in a subgroup isomorphic to $A_4$, $D_{10}$, or
$S_3$, where $D_{10}$ has order ten. Indeed, in the natural action on five
letters an intransitive subgroup fixes a letter or has orbits of sizes
$2+3$. A proper transitive subgroup has order $5,10,15,20$, or $30$.
Orders $15$ and $20$ would force a normal Sylow $5$-subgroup, whose
normalizer in $A_5$ has order ten; order $30$ is excluded by simplicity.
The remaining cases lie in a Sylow $5$-normalizer.

Each of $A_4,D_{10},S_3$ has independence invariant two. For $A_4$, any
generating set contains a $3$-cycle and an element outside its cyclic
subgroup; these generate $A_4$, since it has no subgroup of order six.
The proper subgroups of $A_4$ have independence invariant at most two.
For a dihedral group of order $2q$ with $q$ prime, a generating set
contains either a nontrivial rotation and a reflection, or two distinct
reflections; two such elements suffice. Deleting one element from an
independent set in $A_5$ therefore gives
\begin{equation}\label{eq:A5}
 \mup(A_5)\le3.
\end{equation}
The natural degree-five action has a minimal base of size three, so equality
holds. These facts also follow directly from the usual subgroup structure of
$A_5$; simplicity follows from its nonidentity conjugacy-class sizes
$15,20,12,12$.

\begin{lemma}\label{lem:H}
One has $\mup(H)=b(H)=3$ and $\bfth(H)=2$.
\end{lemma}
\begin{proof}
The only involution in $\SL_2(29)$ is $-I$: a matrix squaring to $I$
is diagonalizable with eigenvalues $\pm1$, and determinant one forces
the signs to agree. Every even-order subgroup of $H$ therefore contains
$Z$. Every odd-order subgroup injects into $A_5$ and is trivial or cyclic
of order three or five by the preceding subgroup description.

Let $S$ be independent in $H$. If its images in $H/Z$, as an indexed
family, are independent, then $|S|\le3$. Otherwise some $s\in S$ lies
in $KZ\setminus K$, where $K=\langle S\setminus\{s\}\rangle$.
Thus $Z\nsubseteq K$, so $K$ is trivial or cyclic of odd prime order.
Independence then gives $|S|\le2$. Hence $\mup(H)\le3$.
The degree-five action of $H/Z$ pulls back to an action of $H$ with a
minimal base of size three. Lemma~\ref{lem:breadth} proves
$b(H)=\mup(H)=3$.

A subgroup family with total intersection $1$ must have an odd-order
member. If that member is trivial, a minimal such family has size one.
If it has prime order, some other member intersects it trivially, so
a minimal family has size at most two. Subgroups of orders three and five
give a faithful family of size two. This proves $\bfth(H)=2$.
\end{proof}

We will need an explicit independent triple in $H$. Put
\begin{equation}\label{eq:R}
 R_1=A,\qquad
 R_2=\begin{pmatrix}12&24\\0&17\end{pmatrix},\qquad
 R_3=\begin{pmatrix}25&3\\4&4\end{pmatrix}.
\end{equation}
The identities $R_2=BAB^2$ and $R_3=AR_2(AB)^2$ put these matrices in
$H$. Direct multiplication gives $R_i^2=-I$ and
\[
 \operatorname{ord}(R_2R_3)=4,\qquad
 \operatorname{ord}(R_1R_3)=3,\qquad
 \operatorname{ord}(R_1R_2)=10.
\]
In $H/Z$ the corresponding products have orders $2,3,5$.
Each pair of projective involutions generates a dihedral group; lifting
back to $H$, the pair subgroups have orders
\begin{equation}\label{eq:pairs}
 |\langle R_2,R_3\rangle|=8,\quad
 |\langle R_1,R_3\rangle|=12,\quad
 |\langle R_1,R_2\rangle|=20.
\end{equation}
The subgroup generated by all three has order divisible by
$\operatorname{lcm}(8,12,20)=120$, so it is $H$. Every pair subgroup
is proper, proving independence. Moreover, the first two pair subgroups
intersect in $\langle R_3\rangle$: this subgroup has order four, which
is also $\gcd(8,12)$. The third does not contain $R_3$ and contains $Z$.
Thus the intersection of the three pair subgroups is exactly $Z$.

\section{Subgroups missing the translations}
\begin{lemma}\label{lem:line}
Let $p$ be an odd prime and let $C\le\F_p^\times$ be a cyclic
$2$-group. For the faithful scalar semidirect product
$L=\F_p\rtimes C$, one has $\mup(L)\le2$.
\end{lemma}
\begin{proof}
Put $P=(\F_p,+)$ and consider a subgroup $J\le L$.
If $J\cap P=1$, then $J$ embeds in $C$ and has independence invariant
at most one: a subgroup of a cyclic prime-power group is generated by
an element of maximum order in any generating set.
Otherwise $P\le J$ and $J=P\rtimes D$ for some $D\le C$.
For $D=1$ the assertion is immediate. If $D\ne1$, every generating
set of $J$ contains an element $x=(a,\lambda)$ whose image generates
$D$, because $D$ has a unique maximal subgroup. Writing $m=|D|$, we have
\[
 x^m=\left(a\sum_{j=0}^{m-1}\lambda^j,1\right)=1.
\]
Thus $\langle x\rangle$ is a complement of order $m$.
Some other generator $y$ lies outside it. The subgroup $\langle x,y\rangle$
projects onto $D$; if it intersected $P$ trivially, its order would be
$m$ and it would equal $\langle x\rangle$, a contradiction.
Hence $\langle x,y\rangle=J$. Every independent generating set of
every subgroup of $L$ has size at most two.
\end{proof}

\begin{lemma}\label{lem:stabilizers}
Every line stabilizer in $H$ is cyclic of order dividing four.
In particular, $H$ acts irreducibly on $V=\F_{29}^2$.
\end{lemma}
\begin{proof}
For a line $W$, the scalar action gives a homomorphism
$H_W\to\F_{29}^\times$. In a basis adapted to $W$, a kernel element
has the form $\left(\begin{smallmatrix}1&u\\0&1\end{smallmatrix}\right)$.
If $u\ne0$ it has order $29$, impossible in a group of order $120$.
The homomorphism is therefore injective, and
\[
 |H_W|\mid\gcd(120,28)=4.
\]
Since $|H|>4$, the whole group cannot stabilize a line.
\end{proof}

Now form $G=V\rtimes H$, with multiplication
$(v,h)(w,k)=(v+hw,hk)$, and write $\pi:G\to H$ for projection.

\begin{lemma}\label{lem:properV}
If $K\le G$ does not contain $V$, then $\bfth(K)\le2$.
More precisely, if $K\cap V=1$, then $\mup(K)\le3$;
if $K\cap V$ is a line, then $\mup(K)\le2$.
\end{lemma}
\begin{proof}
Every additive subgroup of $V$ is an $\F_{29}$-subspace.
If $K\cap V=1$, projection embeds $K$ into $H$, so
Lemma~\ref{lem:H} and subgroup monotonicity give both bounds.
Otherwise $W=K\cap V$ is a line, normal in $K$, and
$\pi(K)\le H_W$ is cyclic of order dividing four.
A Sylow $2$-subgroup of $K$ maps isomorphically onto $\pi(K)$
and complements $W$. Its action on $W$ is faithful by
Lemma~\ref{lem:stabilizers}. Lemma~\ref{lem:line} gives
$\mup(K)\le2$, hence also $\bfth(K)\le2$.
\end{proof}

\section{The invariants of the affine group}
\begin{proposition}\label{prop:upper}
One has $\bfth(G)\le b(G)\le3$.
\end{proposition}
\begin{proof}
First let $K_1,\ldots,K_t$ be a meet-irredundant family with total
intersection $1$. Some $K_i$ does not contain $V$.
Lemmas~\ref{lem:restriction} and \ref{lem:properV} imply
\[
 t-1\le\bfth(K_i)\le2,
\]
so $\bfth(G)\le3$.

Every nontrivial normal subgroup $N\lhd G$ contains $V$.
Indeed, irreducibility gives $N\cap V=1$ or $V$.
In the first case $[N,V]\le N\cap V=1$, whence
$N\le C_G(V)=V$, since the matrix action is faithful.
This forces $N=1$, a contradiction. Every nonfaithful action of $G$
therefore factors through $G/V\cong H$, and its minimal bases have
size at most $b(H)=3$.
\end{proof}

Let $e_1=(1,0)^{\mathsf T}$ and $e_2=(0,1)^{\mathsf T}$.
The set
\begin{equation}\label{eq:S}
 S=\{(e_1,I),(0,R_1),(0,R_2),(0,R_3)\}
\end{equation}
is independent: projecting to $H$ prevents any one of the three linear
elements from being generated by the others, and the translation is
outside the zero-translation complement. It also generates $G$, since
the $H$-orbit of $e_1$ spans the irreducible module $V$.

\begin{proposition}\label{prop:mu}
One has $\mup(G)=4$.
\end{proposition}
\begin{proof}
The set \eqref{eq:S} gives the lower bound. For the upper bound consider
a meet-irredundant family $K_1,\ldots,K_t$, without any condition on its
total intersection. If all members contain $V$, quotienting by $V$
gives a meet-irredundant family in $H$, so $t\le3$.
Otherwise choose $K_i\not\supseteq V$. Lemmas~\ref{lem:restriction}
and \ref{lem:properV} give $t-1\le\mup(K_i)\le3$.
Apply Lemma~\ref{lem:breadth}.
\end{proof}

There is a faithful minimal base of size three. Define
\[
\begin{split}
 L_1&=\langle(e_1,I),(0,-I)\rangle,\\
 L_2&=\langle(e_2,I),(0,-I)\rangle,\\
 L_3&=\langle(e_1+e_2,I),(2e_1,-I)\rangle.
\end{split}
\]
Each subgroup has order $58$, and their pairwise intersections are
\[
\begin{split}
 L_1\cap L_2&=\{1,(0,-I)\},\\
 L_1\cap L_3&=\{1,(2e_1,-I)\},\\
 L_2\cap L_3&=\{1,(-2e_2,-I)\}.
\end{split}
\]
These follow by comparing their translation lines and their cosets of
elements with linear part $-I$. Their total intersection is $1$.
Proposition~\ref{prop:bases} gives $\bfth(G)\ge3$, realized on
$G/L_1\sqcup G/L_2\sqcup G/L_3$, of degree $5220$.
Since $|G|=29^2\cdot120=100920$, this completes the proof of
Theorem~\ref{thm:main}.

\section{The normality obstruction}
Let $M_i=\langle S\setminus\{s_i\}\rangle$, with the elements of
\eqref{eq:S} in the displayed order. Then $M_1=H$ and the other
three subgroups are the preimages under $\pi$ of the pair subgroups
in \eqref{eq:pairs}. To see the latter assertion, each pair subgroup
has order greater than four and hence cannot stabilize a line by
Lemma~\ref{lem:stabilizers}; its orbit of $e_1$ spans $V$.
Consequently
\[
 \bigcap_{i=1}^4 M_i=Z_0,\qquad
 Z_0=\{(0,I),(0,-I)\}.
\]
This subgroup is not normal in $G$, since
\[
 (e_1,I)(0,-I)(-e_1,I)=(2e_1,-I)\notin Z_0.
\]
Thus $S$ gives a Boolean meet-semilattice of rank four with nonnormal
bottom, whereas Theorem~\ref{thm:main} excludes any such rank-four
embedding with normal bottom.

Taking normal cores destroys irredundancy. The three pair subgroups
have core $Z$ in $H$, by simplicity of $H/Z$. Their preimages therefore
have core $V\rtimes Z$ in $G$. The complement has trivial core, since
every nontrivial normal subgroup of $G$ contains $V$.
The four cores are $1,V\rtimes Z,V\rtimes Z,V\rtimes Z$,
whose intersection is already determined by the first member.

\section*{Reproducibility and use of artificial intelligence}
The proof above is independent of subgroup enumeration. A supplementary
verification package supplies exact Python and C++ implementations that
reconstruct the matrix group, compare all its subgroups, and check the
affine witnesses. It also records a characteristic-eleven example showing
why the line-stabilizer hypothesis matters. The current package is available
in the \href{https://github.com/AlecKriebel/Math/tree/be0f06ee0c61739dcda5629e2c4223e7da7951d3/kourovka_16_45/journal_submission_2026_09_25/supplement}{repository snapshot of 25 September 2026}.
The preprint and original
verification materials are archived at
\href{https://doi.org/10.5281/zenodo.22929486}{doi:10.5281/zenodo.22929486}.

OpenAI Codex (\url{https://openai.com/codex/}) was used during
19--25 September 2026 to develop and check the construction and proofs,
generate manuscript text and exact verification programs, search and check
references, and prepare the submission. Separate AI agents performed
adversarial structural, computational, and source reviews. Historical
application and model build identifiers were not consistently recorded.
These checks are not external peer review or proof-assistant certification.
The author is responsible for the mathematical claims and the submitted text.

\section*{Funding}
This research received no external funding.
The author declares no competing interests.

\begin{thebibliography}{9}
\bibitem{ab2024}
M.~Anagnostopoulou-Merkouri and T.~C. Burness,
\emph{On the regularity number of a finite group and other base-related
invariants}, J. Lond. Math. Soc. \textbf{110} (2024), no.~6, e70035.
\url{https://doi.org/10.1112/jlms.70035}.

\bibitem{cameron2014}
P.~J. Cameron, \emph{Some measures of finite groups related to permutation
bases}, arXiv:1408.0968 (2014).
\url{https://arxiv.org/abs/1408.0968}.

\bibitem{cameron2024}
P.~J. Cameron, \emph{Independence and bases: theme and variations},
Model Theory \textbf{3} (2024), no.~2, 417--431.
\url{https://doi.org/10.2140/mt.2024.3.417}.

\bibitem{notebook}
E.~I. Khukhro and V.~D. Mazurov (eds.),
\emph{Unsolved Problems in Group Theory: The Kourovka Notebook},
21st edition, September 2026 update, Problem 16.45, p.~100.
\href{https://kourovkanotebookorg.wordpress.com/wp-content/uploads/2026/09/21tkt.pdf}{Official notebook and updates}.
\end{thebibliography}
\end{document}
